\documentclass[10pt]{article}
\usepackage[margin=0.9in]{geometry}
\usepackage[T1]{fontenc}
\usepackage{lmodern}
\usepackage{amsmath,amssymb,amsthm,mathtools,microtype}
\usepackage[hidelinks]{hyperref}
\usepackage{enumitem}
\setlist[enumerate]{leftmargin=*,itemsep=2pt,topsep=4pt}
\newtheorem{theorem}{Theorem}
\newtheorem{lemma}{Lemma}
\newtheorem{corollary}{Corollary}
\newcommand{\E}{\mathbb E}
\newcommand{\Prob}{\mathbb P}
\newcommand{\Var}{\operatorname{Var}}
\newcommand{\Poi}{\operatorname{Poi}}
\newcommand{\supp}{\operatorname{supp}}
\newcommand{\eff}{\operatorname{eff}}
\newcommand{\ind}{\mathbf 1}
\DeclareMathOperator{\arcosh}{arcosh}
\DeclareMathOperator{\arsinh}{arsinh}
\DeclareMathOperator{\artanh}{artanh}
\title{An Explicit Constant for Wang's Support-Size Tester}
\author{}
\date{}
\begin{document}
\begin{center}
{\Large\bfseries An Explicit Constant for Wang's Support-Size Tester\par}
\end{center}
\vspace{-0.5em}
\begin{abstract}
Wang's optimal support-size tester uses the explicit sampling coefficient
$C_0=2^{26}$. We show that $C_0=1025$ suffices with every other parameter
and decision rule unchanged, including the branch threshold $L_0=1024$.
The resulting Poissonized tester retains error below $1\%$ on each promise
class and a $98\%$ joint guarantee for its support-size lower bound.
The proof replaces a loose polynomial-coefficient estimate with a
factorial-weight integral bound and sharpens the analysis of the counting
branch. An explicit finite distribution shows that no positive integer
$C_0\le1024$ preserves the $1\%$ guarantee for this unchanged construction.
\end{abstract}

\section{Construction and main result}
Let $p$ be a probability distribution on a finite or countable alphabet.
Write $S=|\supp(p)|\in\mathbb N\cup\{\infty\}$ and
$\Delta_r(p)=\sum_{i>r}p_{(i)}$, where the probabilities are in nonincreasing
order, padded with zeros if necessary. Given an integer $n\ge1$ and
$0<\epsilon<1$, support-size testing distinguishes
$S\le n$ from $\Delta_n(p)\ge\epsilon$. Define
$\eff_\epsilon(p)=\min\{k\ge1:\Delta_k(p)\le\epsilon\}$.
Here $\epsilon$ is the separation parameter, distinct from the target
failure probability $0.01$.
All logarithms are natural unless explicitly subscripted.

We retain the construction of Wang~\cite[Section~3]{wang}, writing $C$ for
its sampling coefficient $C_0$. Set
\begin{equation}
 L=\log(en),\qquad m=\frac{Cn}{\epsilon L},\qquad
 M\sim\Poi(m),\qquad N_i\sim\Poi(mp_i)\ \text{independently},
 \label{eq:budget}
\end{equation}
where $N_i$ counts label $i$ among the $M$ iid draws from $p$, and let
$D=\sum_i\ind\{N_i>0\}$. If $L<1024$, set $\widetilde S=D$.
Otherwise, put
\begin{align}
 a_0&=\frac{\epsilon}{8n},& b&=\frac{64L}{m},&
 J&=\left\lfloor\log_2\frac b{2a_0}\right\rfloor,\label{eq:scales1}\\
 a_j&=2^ja_0,& \delta_j&=2^{-j-2},&
 k_j&=\left\lceil\sqrt{b/a_j}\log(2/\delta_j)\right\rceil
 \quad(0\le j\le J).\label{eq:scales2}
\end{align}
Let $T_k$ be defined by $T_k(\cos\theta)=\cos(k\theta)$.
For generic parameters $(a,b,k,m)$, define the polynomial and weights
\begin{align}
 P(x)&=-\rho T_k\!\left(\frac{b+a-2x}{b-a}\right)
       =\sum_{\ell=0}^k c_\ell x^\ell,
 &\rho&=T_k\!\left(\frac{b+a}{b-a}\right)^{-1},\label{eq:polynomial}\\
 f(0)&=0,\qquad
 f(\ell)=1+c_\ell\ell!/m^\ell\quad(1\le\ell\le k),
 & f(\ell)&=1\quad(\ell>k).\label{eq:weights}
\end{align}
Use subscripts $j$ for $(a_j,b,k_j,m)$, and set
\begin{equation}
 Z_j=\sum_i f_j(N_i),\qquad
 \widetilde S=\max\left\{D,\max_{0\le j\le J}
                  \frac{Z_j-\delta_jn}{1+\delta_j}\right\}.
 \label{eq:statistic}
\end{equation}
The tester returns Yes iff $\widetilde S\le n$. Its lower-bound output is
$\widehat S=\min\{\max\{\lceil\widetilde S\rceil,1\},n\}$.

\begin{theorem}\label{thm:main}
For $C=1025$, the construction above satisfies, uniformly over $p,n,\epsilon$,
\begin{align}
 S\le n&\quad\Longrightarrow\quad
     \Prob(\widetilde S>S)<0.01,\label{eq:upper}\\
 1\le r\le n,\quad\Delta_r(p)\ge\epsilon
     &\quad\Longrightarrow\quad
     \Prob(\widetilde S\le r)<0.01.\label{eq:lower}
\end{align}
Thus the tester has error below $1\%$ on each promise class, and
\begin{equation}
 \Prob\!\left\{\min\{\eff_\epsilon(p),n\}\le\widehat S\le S\right\}
 \ge0.98.\label{eq:coverage}
\end{equation}
\end{theorem}

Only $C$ changes from the original construction, which uses $2^{26}$;
the reduction in the Poisson mean is $2^{26}/1025>65{,}472$.
Our proof replaces an intermediate weight bound rather than assuming
that the original analysis remains valid after substituting a smaller
constant. Integer optimality for this fixed construction and error target
is established in Theorem~\ref{thm:integer}; it is not a minimax constant
claim over all testers.

\section{Two auxiliary bounds}
\begin{lemma}[Discovery-time domination]\label{lem:count}
If $\Delta_r(p)\ge\epsilon$, then
$\Prob(D\le r)\le\Prob\{\Poi(m\epsilon)\le r\}$.
\end{lemma}
\begin{proof}
Embed sampling in a rate-one marked Poisson process. After $k\le r$
distinct labels have been discovered, their total mass is at most the mass
of the $r$ largest categories; hence the unseen mass $q_k$ is at least
$\epsilon$. Conditional on the discovered labels, the next discovery has
waiting time $E_k/q_k$, where $E_k\sim\mathrm{Exp}(1)$ is fresh and
independent of the next label, drawn from the normalized unseen masses.
This constructs the discovery process using independent $E_0,\ldots,E_r$.
Its first $r+1$ waiting times are bounded by $E_k/\epsilon$, so their sum
is stochastically dominated by $\mathrm{Gamma}(r+1,\epsilon)$.
The latter's survival probability at time $m$ is the asserted Poisson tail.
\end{proof}

\begin{lemma}[Expectation and factorial-weight bounds]\label{lem:weights}
Let $0<a<b\le1$, $m>0$ and an integer $k\ge1$ be given, and define $P,\rho,f$ by
\eqref{eq:polynomial}--\eqref{eq:weights}. If $m\sqrt{b(b-a)}\ge k$, then
\begin{equation}
 (1-\rho)\min\{x/a,1\}
 \le Q(x):=\E f(\Poi(mx))\le1+\rho\qquad(x\ge0).
 \label{eq:expectation}
\end{equation}
Moreover, without the condition $m\sqrt{b(b-a)}\ge k$,
\begin{equation}
 \|f\|_\infty\le4\exp\!\left(\frac{2k^2}{m(b-a)}\right).
 \label{eq:norm}
\end{equation}
\end{lemma}
\begin{proof}
The Poisson series gives $Q(x)=1+e^{-mx}P(x)$. On $[0,a]$, $P$ is
nonpositive, increasing and concave, with $P(0)=-1$ and $P(a)=-\rho$;
its chord yields $(1-\rho)x/a\le Q(x)\le1$. On $[a,b]$, $|P(x)|\le\rho$.
For $x\ge b$, put $u=2(x-b)/(b-a)$. The identities for Chebyshev
polynomials~\cite[Section~18.5]{dlmf} give
$|P(x)|=\rho T_k(1+u)\le\rho e^{k\sqrt{2u}}$.
Since $x^2\ge4b(x-b)$ and $m\sqrt{b(b-a)}\ge k$, we have
$mx\ge k\sqrt{2u}$. Thus $|Q(x)-1|\le\rho$, proving
\eqref{eq:expectation}.

For \eqref{eq:norm}, let $z_0=(b+a)/(b-a)>1$. The roots of $T_k$ lie
in $(-1,1)$, so every coefficient of $T_k(z_0+v)$ in $v$ is positive.
As $c_0=-1$, this gives
\begin{align}
 \|f\|_\infty
 &\le \sum_{\ell=0}^k |c_\ell|\ell!/m^\ell\notag\\
 &=\int_0^\infty e^{-s}
 \frac{T_k(z_0+2s/[m(b-a)])}{T_k(z_0)}\,ds.
 \label{eq:integral}
\end{align}
The equality uses $\int_0^\infty e^{-s}s^\ell ds=\ell!$.
For $v,u\ge0$, the identity
$\arcosh(1+2v)=2\arsinh\sqrt v$ and the $1$-Lipschitz property of
$\arsinh$ imply
\[
 \arcosh(1+2v+2u)-\arcosh(1+2v)\le2\sqrt u.
\]
Using $T_k(z)=\cosh(k\arcosh z)$, the ratio in \eqref{eq:integral} is
at most $2\exp(2k\sqrt{s/[m(b-a)]})$. Finally,
$2k\sqrt{s/[m(b-a)]}\le s/2+2k^2/[m(b-a)]$; integration of $e^{-s/2}$
establishes \eqref{eq:norm}.
\end{proof}

\section{Proof of Theorem~\ref{thm:main}}
\paragraph{The counting branch: $L<1024$.}
For $x>1$, the exponential Markov inequality gives
\begin{equation}
 \Prob\{\Poi(nx)\le n\}\le e^{-ng(x)},\qquad
 g(x)=x-1-\log x,\quad
 g(1+h)\ge\frac{h^2}{2(1+h)}\quad(h\ge0).
 \label{eq:chernoff}
\end{equation}
Here $g$ is increasing on $(1,\infty)$ and the last inequality follows by
integrating $g'(1+h)=h/(1+h)$. If $n<10^8$, then $L<20$ and
$ng(1025/L)\ge g(1025/20)>46$. If $n\ge10^8$, then $1025/L>1+1/1024$,
so
\[
 ng(1025/L)>\frac{10^8}{2\cdot1024\cdot1025}>47.
\]
Both exceed $\log100$. Lemma~\ref{lem:count} and $r\le n$ prove
\eqref{eq:lower}; $D\le S$ proves \eqref{eq:upper} in this branch.

\paragraph{Parameters in the polynomial branch: $L\ge1024$.}
The following bounds hold throughout $1024\le C\le2^{26}$.
Since $n=e^{L-1}$,
\begin{equation}
 \frac b{2a_0}=\frac{256L^2}{C}\ge4,\qquad
 0<a_j\le b/2,\qquad b\le\frac{L^2}{16e^{L-1}}<1.
 \label{eq:parameters}
\end{equation}
The last function is decreasing for $L\ge1024$. Furthermore,
\[
 k_j\le\sqrt{512/C}\,L\frac{(j+3)\log2}{2^{j/2}}+1
 \le\left(\frac{3\log2}{\sqrt2}+\frac1{1024}\right)L<\frac32L,
\]
because $(j+3)2^{-j/2}$ decreases for integers $j\ge0$.
Consequently $m\sqrt{b(b-a_j)}\ge64L/\sqrt2>k_j$ and
$m(b-a_j)\ge32L$. Also
\begin{equation}
 \rho_j=\frac1{\cosh(2k_j\artanh\sqrt{a_j/b})}
 \le2e^{-2k_j\sqrt{a_j/b}}\le\delta_j^2/2\le\delta_j.
 \label{eq:rho}
\end{equation}
Lemma~\ref{lem:weights} therefore yields
\begin{align}
 (1-\delta_j)\sum_i\min\{p_i/a_j,1\}
 &\le\E Z_j\le(1+\delta_j)S,\label{eq:means}\\
 \|f_j\|_\infty&\le4e^{9L/64}.\label{eq:uniformnorm}
\end{align}

\paragraph{Simultaneous concentration.}
Independence and $f_j(0)=0$ give
$\Var(Z_j)\le\|f_j\|_\infty^2\E D$.
This remains valid on a countable alphabet: if a truncation omits mass
$\tau$, its error is bounded in absolute value by $\|f_j\|_\infty M_{\rm tail}$,
where $M_{\rm tail}\sim\Poi(m\tau)$; its second moment tends to zero.
If $\E D\le2n$, Chebyshev's inequality and \eqref{eq:uniformnorm} imply
\[
 \Prob\{|Z_j-\E Z_j|>\delta_jn\}
 \le\frac{32e^{9L/32}}{n\delta_j^2}.
\]
Since $\delta_J^{-1}=4\cdot2^J\le1024L^2/C\le L^2$ and
$\sum_{j=0}^J\delta_j^{-2}\le(4/3)L^4$, the event
$\mathcal E=\bigcap_{j=0}^J\{|Z_j-\E Z_j|\le\delta_jn\}$ satisfies
\begin{equation}
 \Prob(\mathcal E^c)\le\frac{128e}{3}L^4e^{-23L/32}
 \le\frac{128e}{3}\,1024^4e^{-736}<0.01.
 \label{eq:concentration}
\end{equation}
The middle bound follows from $4/L-23/32<0$. The final inequality needs
no delicate numerical approximation: $2<e<3$ bounds its left side by
$2^{-689}<1/100$.

\paragraph{Upper coverage.}
If $S\le n$, then $\E D\le n$ and, on $\mathcal E$,
$(Z_j-\delta_jn)/(1+\delta_j)\le\E Z_j/(1+\delta_j)\le S$.
Together with $D\le S$, this proves \eqref{eq:upper}.

\paragraph{Power.}
Fix $1\le r\le n$ with $\Delta_r(p)\ge\epsilon$. If $\mu=\E D>2n$,
then $\Var(D)\le\mu$ and
\[
 \Prob(\widetilde S\le r)\le\Prob(D\le n)
 \le\frac{\mu}{(\mu-n)^2}\le\frac2n<0.01.
\]
Otherwise $\E D\le2n$. Consider the following exhaustive cases.
\begin{enumerate}[label=(\roman*)]
\item $p_{(r)}\ge b/2$. Observing all top $r$ labels and any tail label
forces $D>r$. Thus
\[
 \Prob(D\le r)\le re^{-mb/2}+e^{-m\Delta_r(p)}
 \le ne^{-32L}+e^{-Cn/L}<0.01.
\]
Indeed, $ne^{-32L}=e^{-31L-1}$ and
$Cn/L\ge1024e^{L-1}/L\ge e^{1023}$.
\item $p_{(r)}<a_0$. Equation~\eqref{eq:means} gives
$\E Z_0\ge(1-\delta_0)\epsilon/a_0=6n$. On $\mathcal E$,
$\widetilde S\ge(6n-2\delta_0n)/(1+\delta_0)=22n/5>r$.
\item $a_0\le p_{(r)}<b/2$. Choose $j\le J$ such that
$a_j\le p_{(r)}<2a_j$. For tail labels,
$\min\{p_i/a_j,1\}\ge p_i/p_{(r)}$, hence
$\E Z_j\ge(1-\delta_j)(r+\epsilon/(2a_j))$. Since
$a_j\delta_j=\epsilon/(32n)$, on $\mathcal E$,
\[
 \widetilde S\ge\frac{\E Z_j-2\delta_jn}{1+\delta_j}
 \ge r+\frac{\delta_j\bigl((14-16\delta_j)n-2r\bigr)}{1+\delta_j}>r,
\]
where $\delta_j\le1/4$ and $r\le n$ make the numerator strictly positive.
\end{enumerate}
These cases prove \eqref{eq:lower}; their error bounds are not added,
since the applicable case is determined by the fixed distribution.

\paragraph{The lower-bound output.}
If $S\ge n$, then $\widehat S\le S$ automatically; otherwise
\eqref{eq:upper} and integrality of $S$ give this with probability at least
$0.99$. Set $r=\min\{\eff_\epsilon(p),n\}-1$. If $r=0$, the other
inequality is automatic. Otherwise $\Delta_r(p)>\epsilon$, so
\eqref{eq:lower} gives $\widehat S\ge r+1$ with probability at least $0.99$.
A union bound proves \eqref{eq:coverage}, completing the proof.
\qed

\begin{corollary}[Rounding and a deterministic sample cap]\label{cor:rounding}
Theorem~\ref{thm:main} remains valid when $m$ is rounded upward and $b$ is
recomputed. A deterministic cap of $\lceil4m\rceil$ samples preserves the
original main theorems' success requirement of at least $2/3$.
\end{corollary}
\begin{proof}
Rounding only increases the count-branch budget. In the polynomial branch,
it replaces $C$ by $C'\in[1025,1025+\epsilon L/n)$; as
$L/n\le1024e^{-1023}<1$, all bounds above still apply.
For a sample cap, draw $M\sim\Poi(m)$ and abort with arbitrary output if
$M>4m$. Chernoff's inequality gives
$\Prob(M>4m)\le e^{-m}\le e^{-1025}$, since $n/\log(en)\ge1$.
The additional failure probability does not affect the $2/3$ requirement.
\end{proof}

\section{Integer optimality for the unchanged construction}
\begin{theorem}\label{thm:integer}
For every $0<C\le1024$, there is a finite distribution with
$\Delta_n(p)=\epsilon$ on which the unchanged tester has error greater
than $0.06$. Consequently $1025$ is its smallest positive integer sampling
coefficient achieving uniform error below $1\%$.
\end{theorem}
\begin{proof}
Let $n=\lceil e^{1023}\rceil-1$. Then $1023<L<1024$, so the tester uses
only $D$. At $C=1024$, write
\begin{equation}
 \frac{1024n}{L}=n+d,\qquad
 0<d=\frac{n\log(e^{1023}/n)}{L}\le\frac1L<\frac1{1023}.
 \label{eq:excess}
\end{equation}
To lower-bound a Poisson tail without an asymptotic approximation, let
$Y\sim\Poi(n)$ and $Z=n-Y$. Its central moments satisfy
$\E Z=0$, $\E Z^2=n$, and $\E Z^4=3n^2+n\le4n^2$.
H\"older and Cauchy--Schwarz give, respectively,
\[
 \E|Z|\ge\frac{(\E Z^2)^{3/2}}{(\E Z^4)^{1/2}}\ge\frac{\sqrt n}{2},
 \qquad
 \Prob(Y<n)\ge\frac{(\E Z_+)^2}{\E Z_+^2}\ge\frac1{16},
\]
where $\E Z_+=\E|Z|/2$ and $\E Z_+^2\le n$. Coupling an independent
$V\sim\Poi(t)$ yields, for every $t\ge0$,
\begin{equation}
 \Prob\{\Poi(n+t)\le n-1\}\ge e^{-t}/16.
 \label{eq:poissonlower}
\end{equation}

Set $\epsilon=1/4$, $K=100n^2+n-1$, and define one head atom and $K$
equal tail atoms by
\[
 q=\frac1{400n^2},\qquad p_0=1-Kq,\qquad p_i=q\ (1\le i\le K).
\]
Here $p_0>q>0$, the masses sum to one, and
$\Delta_n(p)=(K-n+1)q=\epsilon$. The number of tail observations is
$Y_K\sim\Poi(mKq)$, and $D\le Y_K+1$. At $C=1024$,
\[
 mKq=(n+d)\left(1+\frac{n-1}{100n^2}\right)=n+\tau,\qquad
 0<\tau<\frac1{1023}+\frac{1+1/1023}{100}<0.012.
\]
Therefore \eqref{eq:poissonlower} gives
$\Prob(D\le n)\ge e^{-\tau}/16>(1-0.012)/16>0.06$.
For smaller $C$, the same distribution has a smaller tail mean and at
least as large a failure probability. Combine with Theorem~\ref{thm:main}.
\end{proof}

The obstruction also survives rounding $m$ upward: the tail mean then
increases by less than $Kq<0.2525$, giving failure probability greater
than $(1-0.012-0.2525)/16>0.01$. Integer optimality is conditional on the
fixed threshold $L_0=1024$ and the $1\%$ target. It neither identifies the
optimal real coefficient nor precludes smaller constants for a modified
algorithm or a weaker success requirement.

\paragraph{Verification and implementation.}
Independent artifact checks cover exact-rational endpoint inequalities,
polynomial coefficients, finite-distribution probabilities, and final
outputs. They supplement the analytic proof and are not a proof-assistant
certificate. Certified evaluation must escalate interval precision until
the decision and, when applicable, the integer output are resolved.

\begin{thebibliography}{2}
\bibitem{wang}
Qisheng Wang.
\newblock \emph{Optimal Support Size Testing}.
\newblock arXiv preprint arXiv:2610.18723v1, 2026.

\bibitem{dlmf}
NIST Digital Library of Mathematical Functions.
\newblock Chapter 18, Classical Orthogonal Polynomials, Sections 18.5 and 18.16.
\newblock \url{https://dlmf.nist.gov/18.5} and \url{https://dlmf.nist.gov/18.16}.
\end{thebibliography}
\end{document}
